\scanpage{50}
% archive foliation: 49

Donc on peut \deleted{considérer que} \inserted{identifier} les $(C,\alpha)$ \deleted{sont} entre eux \deleted{\unsure{et tout}}, et trouver un
\[
[C(\xi),\alpha(\xi)]\quad\text{canonique}.
\]

\deleted{On peut alors considérer sur $A''$ \inserted{$=A'\times A_\infty$} (anneau des adèles premières à $p$), la gerbe $g''$ \inserted{de groupe $\mu_2$} qui sur $A'$ est la gerbe des racines carrées de $A'(1)$, et sur $A_\infty$ est la gerbe type $g'_\infty$ déjà envisagée \unsure{plus haut}.}

En fonction de ce choix, on trouve une gerbe canonique $\mathfrak G(\xi)$ sur l'anneau des adèles, de groupe $\mu_2$, savoir la gerbe qui sur $A''=A'\times A_\infty$ est la gerbe des équivalences de la gerbe $\unsure{\mathfrak g}_{\mathbb Z[1/p]}$ \inserted{de groupe $\mu_2$} définie sur $\mathbb Z[1/p]$ par $(C,T)$ \deleted{\unsure{avec la}} avec la gerbe $g'_{A''}$ déjà envisagée plus haut (sur $A'$ c'est celle des $\sqrt{\phantom{a}}$ de $A'(1)$, sur $A_\infty$ c'est $g_\infty$ déjà envisagée plus haut) et qui sur $A_p$ est \uline{la gerbe \deleted{de}\inserted{triviale}}.
Cette gerbe est triviale, et l'ens des classes \inserted{sections} à isom \unsure{près} est le \uline{torseur} $\mathcal T$ sous $H^1(\mathbb A,\mu_2)$ déjà \inserted{envisagé plus haut, qui est can. indép. de $\xi$.}
\deleted{\unsure{qui ne} \unsure{dépend pas du} \unsure{couple}}
\deleted{\unsure{définie par} la \uline{section canonique} \unsure{de} $\unsure{\mathfrak G}(\xi)$}
\margin{\deleted{\unsure{Autrement}}}

Tout ceci \unsure{posé|fait}, la théorie \inserted{de Deuring} des courbes elliptiques d'inv. de Hasse nul nous fournit

a) un élément $\xi_p\in\mathcal T_p$ (\uline{l'élément de Deuring})

b) une section \deleted{canonique} $\xi$ de $\mathfrak G(\xi_p)$, définie à isom près, unique \uline{au signe près à l'infini}, dont la valeur en $p$ \deleted{\unsure{canonique}} est la section unité \deleted{\unsure{canonique}} de $\mathfrak G(\xi_p)_{\mathbb Z_p}$, et la valeur en $\infty$ est la \inserted{polarisation donnée par} $\unsure{\Gamma|T}$\ldots

On trouve donc un \deleted{\unsure{signe}}\inserted{élément} $|\xi|$ de $\mathcal T$, \unsure{dont} la compos. suivant $p$ est $\xi_p$.
